% =========================================================================
% STANDALONE COMPILABLE VERSION — tested with pdflatex
% To embed in the main roadmap document:
%   1. Copy all \definecolor + \newcommand declarations into your preamble
%   2. Copy the \begin{landscape}...\end{landscape} block into the body
%   3. Required packages (add to preamble if absent):
%        \usepackage{xcolor}
%        \usepackage{colortbl}
%        \usepackage{array}
%        \usepackage{longtable}
%        \usepackage{pdflscape}
%        \usepackage{geometry}
% =========================================================================

\documentclass[10pt,a4paper]{article}
\usepackage[T1]{fontenc}
\usepackage{geometry}
\usepackage{xcolor}
\usepackage{colortbl}
\usepackage{array}
\usepackage{longtable}
\usepackage{pdflscape}

\geometry{margin=1.5cm}

% ---- Row background colours ----
\definecolor{gaterow}{RGB}{240,240,238}
\definecolor{currentrow}{RGB}{253,243,237}
\definecolor{phase1row}{RGB}{237,245,253}
\definecolor{phase2row}{RGB}{237,249,244}
\definecolor{phase3row}{RGB}{245,244,254}
\definecolor{endstaterow}{RGB}{238,237,254}
\definecolor{colheader}{RGB}{245,245,243}

% ---- SSL badge colours ----
\definecolor{sslredbg}{RGB}{252,235,235}
\definecolor{sslredtx}{RGB}{163,45,45}
\definecolor{sslamberbg}{RGB}{250,238,218}
\definecolor{sslambtx}{RGB}{133,79,11}
\definecolor{ssltransbg}{RGB}{230,241,251}
\definecolor{ssltrantx}{RGB}{24,95,165}
\definecolor{ssltealbg}{RGB}{225,245,238}
\definecolor{ssltealtx}{RGB}{15,110,86}
\definecolor{sslpurbg}{RGB}{238,237,254}
\definecolor{sslpurtx}{RGB}{60,52,137}

% ---- Commands ----
\newcommand{\sslbadge}[3]{\colorbox{#1}{\scriptsize\textbf{\textcolor{#2}{#3}}}}
\newcommand{\kpistrip}[1]{\par\vspace{2pt}\noindent\rule{\linewidth}{0.2pt}\par\noindent\textit{\scriptsize #1}}
\newcommand{\cellhead}[1]{\textbf{#1}\par\smallskip}
\newcommand{\gatelbl}[1]{\scriptsize\textit{Gate~#1}}

\begin{document}
\begin{landscape}

\setlength{\tabcolsep}{5pt}
\renewcommand{\arraystretch}{1.3}

% Landscape A4 with 1.5cm margins: textwidth ~26.7cm
% Columns: 3.0cm label + 4 x 5.6cm data = 25.4cm

\begin{longtable}{%
|>{\footnotesize\raggedright\arraybackslash}p{3.0cm}%
|>{\footnotesize\raggedright\arraybackslash}p{5.6cm}%
|>{\footnotesize\raggedright\arraybackslash}p{5.6cm}%
|>{\footnotesize\raggedright\arraybackslash}p{5.6cm}%
|>{\footnotesize\raggedright\arraybackslash}p{5.6cm}|}

\caption{Master Capability Transition Framework: Phase States, Gate
Conditions, and SSL Bands across Four Analytical Dimensions.
Gate condition rows state named, verifiable triggers for each phase
transition derived from the Critical Enabler cells of the per-priority
progression matrices. SSL bands track Strategic Sovereignty Level per
phase. KPI anchors serve as cross-reference consistency anchors only.}
\label{tab:master_transition}\\

% ---- Column headers (first page) ----
\hline
\rowcolor{colheader}
\textbf{Phase / SSL Band} &
\textbf{R\&I Maturity}
  \newline{\scriptsize TRL stage $\cdot$ research instruments $\cdot$ innovation outputs} &
\textbf{Policy, Regulatory \& Standards}
  \newline{\scriptsize enforcement $\cdot$ certification $\cdot$ standardisation bodies} &
\textbf{Operational Deployment \& Adoption}
  \newline{\scriptsize pilots $\cdot$ integrations $\cdot$ scale $\cdot$ actor categories} &
\textbf{Market Uptake \& Strategic Impact}
  \newline{\scriptsize market share $\cdot$ SME access $\cdot$ competitiveness $\cdot$ sovereignty} \\
\hline
\endfirsthead

% ---- Column headers (continuation pages) ----
\hline
\rowcolor{colheader}
\textbf{Phase / SSL Band} &
\textbf{R\&I Maturity} &
\textbf{Policy, Regulatory \& Standards} &
\textbf{Operational Deployment \& Adoption} &
\textbf{Market Uptake \& Strategic Impact} \\
\hline
\endhead

\hline\multicolumn{5}{r}{\footnotesize\textit{Continued\ldots}}\\
\endfoot

\hline
\endlastfoot

% =================================================================
% ROW I — CURRENT STATUS
% =================================================================
\rowcolor{currentrow}
\textbf{I.~Current Status}
\newline\textit{2025 baseline}
\newline\sslbadge{sslredbg}{sslredtx}{SSL: High Foreign Dependency}
&
\cellhead{Fragmented, academically dominated}
EU R\&I at TRL~3--5 for sovereign alternatives.
Heavy dependency on NIST PQC process, Intel TDX/AMD SEV-SNP TEE
architectures, and US/Chinese foundation models for SecOps.
No coordinated EU-wide applied R\&I programme for sovereign secure silicon
or edge-native confidential computing.
\kpistrip{Baseline: sovereign secure silicon TRL $\leq$4;
no EU-manufactured EUCC `High' secure element in production.}
&
\cellhead{Regulatory stack adopted, not yet enforced}
NIS2 enforcement initiated (Oct.~2024); CRA implementing technical acts
pending; AI Act high-risk provisions not yet operational; EUCS
non-binding; ENISA PQC guidance absent; NIS2 transposition fragmented
across Member States.
\kpistrip{Baseline: no binding EUCS; no published EU PQC migration
guidance; CRA technical acts not yet in force.}
&
\cellhead{No continental sovereign security infrastructure}
TEE deployment limited to foreign hyperscaler implementations; PQC
confined to experimental pilots; Simpl pre-launch; federated SOC at
bilateral PoC stage only; Island Mode not standardised; supply chain
attestation manual in Open~RAN ecosystems.
\kpistrip{Baseline: 0 Common Data Spaces with Simpl security MIMs;
bilateral SOC PoCs only.}
&
\cellhead{Near-zero EU share in critical security components}
EU domestic production of critical secure silicon near 0\%;
US/Asian hyperscalers dominate cloud security tooling; no EU-sovereign
LLM for SecOps; SME CRA/NIS2 compliance costs unaddressed; RISC-V
security products absent from public procurement frameworks.
\kpistrip{Baseline: EU secure silicon critical infrastructure share
$<$5\%; no EU-sovereign LLM for SecOps in deployment.}
\\
\hline

% =================================================================
% GATE I -> II
% =================================================================
\rowcolor{gaterow}
\gatelbl{I $\rightarrow$ II} &
\multicolumn{4}{>{\footnotesize\raggedright\arraybackslash}p{23.0cm}|}{%
\textit{\textbf{R\&I:} Horizon Europe calls for sovereign secure silicon and
lightweight PQC issued; ECCC work programmes activated.
\quad\textbf{Policy:} CRA implementing technical acts in force; NIS2~Art.~21
national guidance published by competent authorities.
\quad\textbf{Operational:} At least one national pilot per priority domain
initiated under regulatory pull.
\quad\textbf{Market:} ECCC-co-funded SME compliance support scheme
operational; IPCEI-CIS initial deployments contracted.}}
\\
\hline

% =================================================================
% ROW II — PHASE 1
% =================================================================
\rowcolor{phase1row}
\textbf{II.~Phase~1}
\newline\textit{2025--2030}
\newline\textit{Pilot Validation \&}
\newline\textit{Applied Maturity}
\newline\sslbadge{sslamberbg}{sslambtx}{SSL: High $\to$ Emerging EU}
&
\cellhead{Applied R\&D (P3,P6,P8,P9) $\cdot$ Fragmented Piloting (P1,P2,P4,P5,P7,P10)}
Horizon Europe calls targeting sovereign secure silicon, lightweight PQC
for 8/16-bit MCUs, Simpl security MIMs, and LLMSecOps. RISC-V secure element
prototypes targeting TRL~6--7 under EU Chips Act co-funding. RF fingerprinting
and CSI-based device authentication validated on EU-manufactured constrained
hardware. ECCC-funded crypto-agility middleware research initiated per
FIPS~203/204/205 adoption.
\kpistrip{KPI: RISC-V prototypes at TRL~6--7; lightweight PQC benchmarked
on ARM Cortex-M4 class hardware.}
&
\cellhead{Regulatory activation creates compliance market pull}
CRA hardware security implementing rules published; NIS2~Art.~21 guidance
issued; ENISA PQC sector-specific migration guidelines by 2026, mapping
FIPS~203/204/205 to NIS2 obligations; EU AI Act Art.~9 compliance frameworks
operational for high-risk security AI; EUCS certification rollout mandating
TEE workload protection; Simpl-Open launched with initial security MIM
definitions; SBOM templates mandated in EU-funded deployments from 2026.
\kpistrip{KPI: ENISA PQC guidelines published and adopted by 5+ Member
States by 2026.}
&
\cellhead{Pilot deployments under regulatory pull across priority domains}
3--5 RISC-V secure element pilots under Horizon Europe; LLM-assisted SOC
pilots in 5+ national cybersecurity centres under AI Act Art.~9; CAMARA API
profiles in 3+ cross-border 5G corridor pilots by 2027; federated SIEMs
with MITRE ATT\&CK integration; Simpl security pilots in 2+ Common Data
Spaces (Manufacturing, Health); Island Mode fallback tested in
NIS2/DORA-compliant edge architectures.
\kpistrip{KPI: 2+ Common Data Spaces piloting Simpl security MIMs;
federated SIEMs active in 3+ Member States.}
&
\cellhead{Sovereignty trajectory established; SME access initiated}
SME compliance support scheme co-funded by ECCC and national agencies;
IPCEI-CIS initial deployments activating bilateral operator interconnection;
EU secure silicon below 10\% but industrial investment trajectory established
under CRA pull; EUCC `High' referenced in public procurement frameworks for
critical infrastructure ICT.
\kpistrip{KPI: EU secure silicon $<$10\% but investment trajectory
established; SME support scheme operational in 15+ Member States.}
\\
\hline

% =================================================================
% GATE II -> III
% =================================================================
\rowcolor{gaterow}
\gatelbl{II $\rightarrow$ III} &
\multicolumn{4}{>{\footnotesize\raggedright\arraybackslash}p{23.0cm}|}{%
\textit{\textbf{R\&I:} RISC-V prototypes validated at TRL~6--7; lightweight PQC
benchmarked on EU-manufactured constrained hardware.
\quad\textbf{Policy:} ENISA PQC guidelines published; EUCS rollout active;
Simpl-Open launched with security MIM definitions.
\quad\textbf{Operational:} Federated SIEM pilots with standardised CTI sharing
active in 3+ Member States; 2+ Common Data Space Simpl pilots validated.
\quad\textbf{Market:} CRA enforcement generating measurable industrial investment
in sovereign hardware alternatives.}}
\\
\hline

% =================================================================
% ROW III — PHASE 2
% =================================================================
\rowcolor{phase2row}
\textbf{III.~Phase~2}
\newline\textit{2030--2035}
\newline\textit{Industrial Scaling,}
\newline\textit{Standardisation \&}
\newline\textit{Cross-Sector}
\newline\textit{Deployment}
\newline\sslbadge{ssltransbg}{ssltrantx}{SSL: Emerging EU $\to$ Full Autonomy}
&
\cellhead{Applied \& industrial R\&D concurrent with standardisation}
Phase~2 is not purely standardisation. Hardware--software co-design for
neuromorphic/PIM architectures (P3); sovereign EU LLM/SLM for SecOps at AI
Act high-risk compliance (P6); RISC-V TEE designs entering volume production
via Chips~JU by 2032 (P1,~P8); FHE performance optimisation programmes
initiated (P8); ZKP-based supply chain attestation R\&D (P5); Cognitive
Twin prototype development (P9).
\kpistrip{KPI: At least one RISC-V TEE design in volume production by 2032;
sovereign EU LLM/SLM certified at AI Act high-risk level.}
&
\cellhead{Standardisation milestones convert pilots into binding frameworks}
ETSI TC~CYBER hybrid PQC profiles published by 2031 (ML-KEM/FIPS~203;
ML-DSA/FIPS~204); CEN/CENELEC JTC~13 Simpl security MIMs ratified by 2032;
3GPP SA3/ETSI TC~ITS Island Mode fallback APIs standardised by 2031;
Runtime SBOM mandate for critical infrastructure supply chains by 2031;
classical-only cryptography deprecated for high-security EU government
contexts by 2033; Data Visa legal instruments piloted.
\kpistrip{KPI: CEN/CENELEC Simpl MIMs ratified; ETSI hybrid PQC profiles
published; classical-only deprecated in sovereign contexts by 2033.}
&
\cellhead{Continental-scale federated deployment across regulated sectors}
First EU RISC-V secure element at EUCC `High' by 2031--2032; hybrid PQC
deployed across Industrial IoT, energy grids, and financial systems;
Pan-European CTI platform operational as LLMSecOps engine by 2032; UEBA
baselines harmonised using STIX/TAXII and MITRE ATT\&CK; multi-party
privacy-preserving AI pipelines demonstrated in health and finance by
2033; AI-driven graceful degradation standardised in O-RAN AI/ML framework.
\kpistrip{KPI: 70\%+ Common Data Spaces on Simpl-based security; APT
mean dwell time reduced by 70\% relative to Phase~1 baseline.}
&
\cellhead{EU industrial sovereignty consolidating; 30--40\% critical component share}
EU Chips Act fabrication producing RISC-V TEE designs in volume; EU secure
silicon at 30--40\% market share in critical infrastructure procurement;
sovereign European edge LLM/SLM in operational SecOps deployment;
IPCEI-CIS outcomes translated into binding inter-operator SLA frameworks
with security audit requirements; Data Visa pilots in Nordic--Baltic and
BeNeLux corridors.
\kpistrip{KPI: EU secure silicon 30--40\% critical infrastructure share;
70\%+ Data Spaces on Simpl security interoperability.}
\\
\hline

% =================================================================
% GATE III -> IV
% =================================================================
\rowcolor{gaterow}
\gatelbl{III $\rightarrow$ IV} &
\multicolumn{4}{>{\footnotesize\raggedright\arraybackslash}p{23.0cm}|}{%
\textit{\textbf{R\&I:} First EU RISC-V TEE design at EUCC `High' and in volume
production; sovereign EU edge LLM/SLM at AI Act high-risk compliance.
\quad\textbf{Policy:} CEN/CENELEC Simpl security MIMs ratified; ETSI hybrid
PQC profiles published; 3GPP Island Mode APIs standardised.
\quad\textbf{Operational:} Pan-European CTI platform operational as LLMSecOps
engine; EU secure silicon exceeding 30\% critical infrastructure share.
\quad\textbf{Market:} EU NaaS consortium offers commercially competitive in
2+ cross-border corridors.}}
\\
\hline

% =================================================================
% ROW IV — PHASE 3
% =================================================================
\rowcolor{phase3row}
\textbf{IV.~Phase~3}
\newline\textit{2035--2040}
\newline\textit{Autonomous Leadership}
\newline\textit{\& Frontier Innovation}
\newline\sslbadge{ssltealbg}{ssltealtx}{SSL: Full Sovereign Autonomy}
&
\cellhead{Frontier research defining global state of the art}
FHE reaching operational edge latency thresholds; zero-energy air interface
security with ambient RF harvesting; self-evolving AI swarm protocols for
autonomous MTD; Cognitive Twin E2E orchestration for cyber-physical
resilience; bio-inspired elasticity protocols standardised via ETSI ENI;
battery-less node quantum-resistant authentication. EuroHPC hosts federated
training infrastructure for self-improving sovereign security AI, preventing
regression to non-EU foundation model dependency.
\kpistrip{KPI: FHE latency within $2\times$ AES-GCM baseline for edge
analytics; EuroHPC sovereign security AI model in operational deployment.}
&
\cellhead{EU transitions from standard-taker to global standard-setter}
EU Cloud Rulebook binding with machine-readable compliance APIs by 2037--2038;
EUCC `High' as default procurement requirement for all new EU ICT products;
PQC-by-default mandated for all new EU ICT products by 2036; AI Act
explainability requirements enforced for autonomous MTD decisions;
``Schengen for Cloud'' legal framework enacted; ETSI/CEN/CENELEC PQC
standards promoted to ISO/ITU-T as global reference implementations.
\kpistrip{KPI: EU Cloud Rulebook binding by 2038; PQC-by-default mandate
in force by 2036; Schengen for Cloud enacted.}
&
\cellhead{Fully autonomous, continent-wide sovereign security operations}
EU Chips Act sovereign fabrication meeting 80\%+ of European secure silicon
demand by 2039; FHE operationalised for real-time edge analytics;
end-to-end confidential continuum from battery-less IoT to HPC;
Cognitive Twins autonomously coordinating E2E cyber-physical recovery;
AI-driven MTD continuously shuffling network resources without human
intervention; SMPC standardised for cross-border federated AI in regulated
sectors (health, finance, defence, justice).
\kpistrip{KPI: EU secure silicon self-sufficiency $\geq$80\% by 2039;
autonomous MTD and Cognitive Twin operational continent-wide.}
&
\cellhead{Competitive parity achieved; standards exported globally}
European operators achieving competitive scale parity with US/Chinese
hyperscalers in cloud-edge NaaS; EU-designed security components dominant
in European market; EUCS/EUCC adopted as global reference certification
frameworks; ``Schengen for Cloud'' eliminating national fragmentation;
EU cybersecurity technology export balance turning positive; SME
participation economically viable without discriminatory compliance barriers.
\kpistrip{KPI: EU NaaS competitive parity in 3+ cross-border corridors;
EUCS/EUCC referenced in ISO/ITU-T standards.}
\\
\hline

% =================================================================
% GATE IV -> V
% =================================================================
\rowcolor{gaterow}
\gatelbl{IV $\rightarrow$ V} &
\multicolumn{4}{>{\footnotesize\raggedright\arraybackslash}p{23.0cm}|}{%
\textit{\textbf{R\&I:} FHE at operational edge latency; zero-energy authentication
validated on battery-less EU hardware; AI swarm protocols stable and
AI Act-explainable.
\quad\textbf{Policy:} EU Cloud Rulebook binding; Schengen for Cloud enacted;
PQC-by-default mandate in force; EU standards submitted to ISO/ITU-T.
\quad\textbf{Operational:} EU secure silicon self-sufficiency $\geq$80\%;
autonomous MTD and Cognitive Twin E2E orchestration operational continent-wide.
\quad\textbf{Market:} EU NaaS competitive parity demonstrated; technology
export balance positive.}}
\\
\hline

% =================================================================
% ROW V — STRATEGIC END-STATE
% =================================================================
\rowcolor{endstaterow}
\textbf{V.~Strategic End-State Objective}
\newline\textit{2040 Vision}
\newline\sslbadge{sslpurbg}{sslpurtx}{SSL: Full Autonomy + Export Leadership}
&
\cellhead{Europe leads the global security R\&I agenda}
EU-originated R\&I outputs --- PQC, TEE architectures, LLMSecOps,
zero-energy security, Cognitive Twins --- define the global research agenda
and are adopted internationally. Aligned with SNS~JU SRIA Phase~3 outcomes:
6G security primitives sovereign and globally deployed. ECCC Strategic Agenda
fully operationalised; European industrial cybersecurity competitiveness
self-sustaining without emergency public subsidy.
\kpistrip{Triple accountability: SNS~JU SRIA $\cap$ ECCC Strategic
Agenda $\cap$ EU regulatory stack.}
&
\cellhead{Complete, coherent, machine-enforceable EU regulatory stack}
NIS2, CRA, AI Act, DORA, Data Act, Cloud Rulebook, and ``Schengen for
Cloud'' operating as a coherent system with machine-readable compliance
APIs and no manual audit requirement. EUCS and EUCC recognised globally
as reference certification frameworks. ETSI and CEN/CENELEC leading
ISO/ITU-T security standardisation tracks.
\kpistrip{End-state: no compliance gap between deployed sovereign
capability and legal obligation; no manual audit required.}
&
\cellhead{Security-by-Design is the default operational baseline}
Autonomous, self-healing, continent-wide security infrastructure requiring
no manual intervention for routine threat response. Full confidential
computing continuum from IoT sensor to HPC node. Zero critical operational
dependencies on non-EU security infrastructure. Island Mode and graceful
degradation are normative operating assumptions, not emergency exceptions.
\kpistrip{End-state: zero critical non-EU security dependencies;
Security-by-Design is default, not option.}
&
\cellhead{Europe: sovereign, competitive, and globally influential}
Europe is the global reference for sovereign, trustworthy, and sustainable
computing security. EU-designed components dominant in European market and
competitive globally. SME participation fully economically viable without
discriminatory compliance costs. European security standards and AI-driven
technologies exported as global reference implementations. Strategic
autonomy sustained without emergency policy intervention.
\kpistrip{End-state: EU cybersecurity technology export balance positive;
SME compliance cost non-discriminatory across all 27~Member States.}
\\
\hline

\end{longtable}
\end{landscape}
\end{document}